\chapter{Model, Score Identities, and Research Context}
\label{ch:background}\label{ch:statmech}\label{ch:stochproc}\label{ch:sde}\label{ch:diffusion}\label{ch:ebm}

This chapter contains everything the research chapters depend on: the data
model (a stationary Gaussian, later Laplace, AR(1) chain), the corruption
model (the Ornstein--Uhlenbeck channel), the identity that converts score
computation into Bayesian inference, the strands of literature the thesis
builds on, and the preliminary work that shaped the final research
questions. Background appears only where a later definition, derivation, or
experiment uses it; for the surrounding theory we cite standard references.

\section{The data model: a stationary AR(1) chain}\label{sec:ar1}

A sequence $(X_0, \dots, X_{K-1})$ is a \emph{Markov chain} if each variable
depends on the past only through its immediate predecessor; its joint
density then factorises into local terms,
\begin{equation}
  p(x_0, \dots, x_{K-1})
  \;=\; p_0(x_0)\, \prod_{k=0}^{K-2} M(x_{k+1} \,|\, x_k),
  \label{eq:chain-factorisation}
\end{equation}
with $M$ the transition kernel \citep{robert2004monte}. The factorisation
\eqref{eq:chain-factorisation} is the structural fact the whole thesis
exploits: it makes the precision matrix of a Gaussian chain tridiagonal
(Chapter~\ref{ch:gaussian}) and the factor graph of the posterior a tree,
on which belief propagation is exact.

The specific chain we study is the first-order autoregressive process, the
canonical Gaussian Markov chain \citep{brockwell1991time}. For
$\alpha \in (-1, 1)$ and i.i.d.\ innovations
$\eta_k \sim \Nd(0, \sigma_\eta^2)$,
\begin{equation}
  a_{k+1} \;=\; \alpha\, a_k + \eta_k ,
  \qquad k = 0, 1, \dots, K-2 ,
  \label{eq:ar1}
\end{equation}
so $M(a' | a) = \Nd(a'; \alpha a, \sigma_\eta^2)$. Throughout we use the
stationary normalisation $\sigma_\eta^2 = 1 - \alpha^2$ and start the chain
in its stationary law, $a_0 \sim \Nd(0,1)$; every frame then has unit
marginal variance. A trajectory $a = (a_0, \dots, a_{K-1})$ is \emph{one}
data point, the prototype of a video or a time series. The frame index $k$
is an internal time, to be kept distinct from the diffusion time $t$
introduced next.

\begin{toolbox}[tb:ar1]{Stationary variance and autocovariance of AR(1)}
\emph{(i) Stationary variance.} If $a_k \sim \Nd(0, \sigma^2)$, then by
independence of $\eta_k$,
$\operatorname{Var}(a_{k+1}) = \alpha^2 \sigma^2 + \sigma_\eta^2$. Imposing
equality of the two variances gives
$\sigma_\infty^2 = \sigma_\eta^2 / (1 - \alpha^2)$, which equals $1$ under
the stationary normalisation. The map
$\sigma^2 \mapsto \alpha^2\sigma^2 + \sigma_\eta^2$ is a contraction for
$|\alpha| < 1$, so any initial variance converges to $\sigma_\infty^2$
geometrically, and Gaussianity is preserved by the affine recursion.

\emph{(ii) Autocovariance.} Unrolling \eqref{eq:ar1} $d$ times,
$a_{k+d} = \alpha^d a_k + \sum_{j=0}^{d-1} \alpha^{j}\eta_{k+d-1-j}$, and
all innovation terms are independent of $a_k$, so in the stationary regime
\begin{equation*}
  \operatorname{Cov}(a_i, a_j) \;=\; \alpha^{|i-j|} .
\end{equation*}
Correlations decay geometrically with the lag, and the covariance matrix of
$K$ consecutive frames is symmetric Toeplitz: constant along diagonals,
because the process is stationary. \hfill$\square$
\end{toolbox}

\section{The corruption model: the Ornstein--Uhlenbeck channel}
\label{sec:ou}

Diffusion models corrupt data with a fixed noising process. The standard
choice \citep{song2021scorebased} is the
Ornstein--Uhlenbeck (OU) process, Brownian motion pulled back to the
origin by a linear spring,
$\dd X_t = -X_t\,\dd t + \sqrt{2}\,\dd W_t$
\citep{uhlenbeck1930theory}. Its transition kernel is Gaussian and
available in closed form \citep{gardiner2009stochastic}: corrupting a
value $a$ for a time $t$ produces
\begin{equation}
  x \;=\; a\, e^{-t} \;+\; \sqrt{\Deltat}\;\varepsilon,
  \qquad \varepsilon \sim \Nd(0, 1),
  \qquad
  \boxed{\;\Deltat \;=\; 1 - e^{-2t}\;} .
  \label{eq:ou-channel}
\end{equation}
The signal is contracted by $e^{-t}$ while isotropic noise of variance
$\Deltat$ fills in: at $t = 0$ the data are untouched; as $t \to \infty$
the output is $\Nd(0,1)$ regardless of the input. The signal-to-noise ratio
$e^{-2t}/\Deltat$ sweeps from $\infty$ to $0$, and it diverges as $1/2t$
for small $t$, a factor that reappears throughout as an error amplifier.
Two consequences are used constantly. First, read as a function of the
\emph{clean} variable, the kernel is again Gaussian,
$\Nd(x; a e^{-t}, \Deltat) \propto
\exp[-(e^{-2t}/2\Deltat)(a - x e^{t})^2]$: a Gaussian likelihood factor of
precision $e^{-2t}/\Deltat$ centred at the deconvolved observation
$e^t x$. Second, pushing a Gaussian vector
$a \sim \Nd(0, \Sigma_0)$ through \eqref{eq:ou-channel} coordinatewise with
independent noises gives a Gaussian output with
\begin{equation}
  \Sigt \;=\; e^{-2t}\, \Sigma_0 \;+\; \Deltat\, I ,
  \label{eq:ou-cov-map}
\end{equation}
since means and covariances transform linearly and the added noise is
white. Equation \eqref{eq:ou-cov-map} is the entire forward process of this
thesis in one line.

\section{Diffusion models and the score}\label{sec:diffusion-idea}
\label{sec:reversal}

Applying the channel \eqref{eq:ou-channel} to data $a \sim p_0$ produces
the noised marginals
\begin{equation}
  p_t(x) \;=\; \int p_0(a)\; \Nd\!\big(x;\, a e^{-t},\, \Deltat I\big)\,
  \dd a ,
  \label{eq:pt-convolution}
\end{equation}
which interpolate from the data distribution to $\Nd(0, I)$. By Anderson's
time-reversal theorem \citep{anderson1982reverse}, the corruption is undone
by a reverse-time SDE whose drift contains a single unknown object, the
\emph{score} of the noised marginals:
\begin{equation}
  \dd X_t \;=\;
  \Big[ -X_t \;-\; 2\, \nabla_x \log p_t(X_t) \Big]\,\dd t
  \;+\; \sqrt{2}\, \dd \bar W_t ,
  \qquad
  \score(x,t) := \nabla_x \log p_t(x) .
  \label{eq:reverse-sde}
\end{equation}
Whoever knows $\score(x,t)$ for all $t$ can sample from $p_0$; the same
score also drives a deterministic transport, the probability-flow ODE
\citep{song2021scorebased}, which connects diffusion models to normalizing
flows and flow matching \citep{chen2018neural, lipman2023flow,
albergo2023building}. In practice the score is approximated by a neural
network trained by denoising score matching
\citep{hyvarinen2005estimation, vincent2011connection, ho2020denoising}:
one regresses a network $s_\theta(x, t)$ onto the kernel score
$\nabla_x \log \Nd(x; a e^{-t}, \Deltat I) = -(x - a e^{-t})/\Deltat$ over
pairs $(a, x)$, which is equivalent to predicting the added noise. The
score is also what makes the energy-based view of generative modelling
\citep{lecun2006tutorial} tractable: for a model
$p_\theta \propto e^{-E_\theta}$ the partition function obstructs
likelihood, sampling, and learning alike, but
$\nabla_x \log p_\theta = -\nabla_x E_\theta$ is free of it. Diffusion
models work entirely with scores and thereby sidestep normalisation.

\section{The score as a posterior expectation}\label{sec:tweedie}

One identity converts score computation into Bayesian inference, and every
result of this thesis passes through it. Differentiating
\eqref{eq:pt-convolution} under the integral sign,
\begin{equation*}
  \nabla_x \log p_t(x)
  = \frac{\int p_0(a)\, \nabla_x \Nd(x; a e^{-t}, \Deltat I)\, \dd a}{p_t(x)}
  = \int p(a \,|\, x)\;
    \frac{a e^{-t} - x}{\Deltat}\; \dd a ,
\end{equation*}
where $p(a|x) = p_0(a)\,\Nd(x; ae^{-t}, \Deltat I)/p_t(x)$ is the posterior
of the clean data given the noisy observation. Hence
\begin{equation}
  \boxed{\;
  \score(x, t)
  \;=\; \frac{e^{-t}\,\E[\,a \,|\, x\,] - x}{\Deltat} .
  \;}
  \label{eq:tweedie}
\end{equation}
We call \eqref{eq:tweedie} the \emph{score--posterior identity}; in
empirical Bayes it is known as Tweedie's formula
\citep{robbins1956empirical, miyasawa1961empirical, efron2011tweedie}. It
holds for \emph{any} data distribution $p_0$: the score of the noisy
marginal and the Bayes-optimal denoiser $\E[a|x]$ determine each other
exactly. Estimating a score is therefore solving an inference problem, and
the quality of any approximate score is the quality of an approximate
posterior mean, amplified by $e^{-t}/\Deltat$.

For sequence data $a \in \R^K$ noised coordinatewise, the identity holds
componentwise with the \emph{full} posterior:
\begin{equation}
  \score_k(x, t)
  \;=\; \frac{e^{-t}\, \E[\,a_k \,|\, x_0, \dots, x_{K-1}\,] - x_k}{\Deltat} .
  \label{eq:tweedie-joint}
\end{equation}
The conditioning is on all frames: although the noise is independent across
frames, the prior couples them, so the $k$-th score component depends on
every observation. This \emph{joint} score is a different object from the
per-frame marginal score $\partial_{x_k}\log p_t(x_k)$, and only the joint
score drives the reverse dynamics \eqref{eq:reverse-sde} of the sequence as
a whole. For our stationary model the distinction is extreme: the marginal
score is degenerate (Section~\ref{sec:development}), while the joint score
carries all the structure. When the prior is a Markov chain, the posterior
in \eqref{eq:tweedie-joint} is a smoothing distribution on a chain, which
is what makes exact computation possible.

\section{Related work}\label{sec:related}\label{sec:cascade}\label{sec:gap}

\paragraph{Statistical physics of learning and of diffusion.} The
statistical-mechanics analysis of neural computation runs from spin
glasses through the Hopfield network and the Boltzmann machine
\citep{edwards1975theory, hopfield1982neural, amit1985storing,
ackley1985learning, smolensky1986information, hinton2002training,
gardner1988space, mezard2017meanfield}, a lineage recognised by the 2024
Nobel Prize in Physics \citep{nobel2024physics}; reviews are
\citet{zdeborova2016statistical, carbone2025hitchhiker}. Its modern branch
studies diffusion models with exact scores: for solvable data
distributions the reverse dynamics crosses sharp regimes, with a
speciation time at which trajectories commit to a mode and, for scores
learned from $n$ samples, a memorisation time below which generation
collapses onto training points \citep{biroli2023generative,
biroli2024dynamical, raya2023spontaneous, ghio2024sampling}. On the
training side, \citet{bachtis2024cascade} show that learning in restricted
Boltzmann machines proceeds through a cascade of second-order phase
transitions: tracking the singular modes of the weight matrix as order
parameters, each learned feature corresponds to a Curie--Weiss critical
point, the critical times are ordered by the spectrum of the data, and the
mean-field signatures (diverging susceptibilities, finite-size scaling,
hysteresis) are observed on MNIST, CelebA, and genomic data
\citep[building on][]{saxe2014exact, decelle2017spectral,
tubiana2017emergence}. Phase transitions thus appear on both the sampling
and the training side of generative modelling. What all of these analyses
hold fixed is the \emph{internal structure of a single data point}: their
coordinates are exchangeable or conditionally i.i.d.; there is no ordering
and no dynamics inside one sample.

\paragraph{Inference on chains.} Exact posterior inference on Gauss--Markov
chains is classical: the Kalman filter and smoother
\citep{kalman1960new, rauch1965maximum}, hidden Markov models
\citep{rabiner1989tutorial}, and their unifying formulation as sum--product
message passing on factor graphs \citep{pearl1988probabilistic,
kschischang2001factor, koller2009probabilistic, mezard2009information,
bishop2006pattern}. This literature fixes the observation noise and asks
for the posterior; it does not ask how the posterior, and with it the
score \eqref{eq:tweedie-joint}, deforms as a function of the diffusion
time $t$.

\paragraph{Message passing with continuous variables.} On trees, belief
propagation computes exact marginals; for continuous variables its
messages are densities, and practical algorithms restrict them to a
parametric family, almost always Gaussian. The classical justification of
Gaussian messages, in approximate message passing and the
Thouless--Anderson--Palmer equations \citep{thouless1977solution,
donoho2009message, zdeborova2016statistical, genovese2026amp}, is a
central-limit argument valid when every node aggregates many weak
incoming messages. On a chain each node has two neighbours, so that
argument does not apply, and we are not aware of a quantitative study of
the Gaussian message closure on chain-structured diffusion posteriors.

\paragraph{The gap.} Sequential data have exactly the structure the
solvable-diffusion programme leaves out: an internal time with Markov
dependence along it. For such data the score is a smoothing posterior on a
chain. How its coupling structure evolves with diffusion time, how it can
be computed by message passing, and what a Gaussian message family costs
when the data are not Gaussian, are the questions addressed in
Chapters~\ref{ch:gaussian} and~\ref{ch:laplace}.

\section{How the research questions took shape}\label{sec:development}

The final model and methods were reached through preliminary work whose
failures were instructive, and we record the path in compressed form.

\emph{Initial question and toy models.} The project began by asking how
diffusion scores behave for data with an internal time, using small
interactive toy models (archived in the project repository): a single
distribution transported by the OU flow; a two-frame additive process
$x_1 = x_0 + c + \eta$ with a multimodal prior and its two-dimensional
joint score; deterministic dynamics (rotations, oscillations) observed
through noise. These experiments made the geometry of score fields
concrete but exposed a conceptual error: scores were being computed
\emph{per frame}, treating a trajectory as $K$ independent samples.

\emph{The joint-score correction.} Supervision corrected the object of
study: a sequence is one sample from a joint law, and the quantity that
drives reverse diffusion is the joint score of the whole noisy trajectory.
For a stationary chain under the variance-preserving channel the
correction is not cosmetic. Each noisy frame is marginally $\Nd(0,1)$ at
every $t$, so the per-frame marginal score is $-x_k$, independent of the
dynamics; all information about the temporal structure lives in the joint
object (Chapter~\ref{ch:gaussian} exhibits this explicitly at $K=2$).

\emph{The minimal model.} This reframing selected the stationary Gaussian
AR(1) chain under coordinatewise OU noise as the minimal model in which
the joint score is nontrivial yet exactly computable, with a single
coupling parameter $\alpha$ against a single noise parameter $t$.

\emph{The message-passing programme and its bottleneck.} Since the score
is a smoothing posterior mean \eqref{eq:tweedie-joint}, supervision
proposed computing it by belief propagation on the chain's factor graph
and asking three questions: whether the recursion reproduces the exact
Gaussian score, what happens for non-Gaussian innovations, and how
accurate local truncations are. Working out the recursion identified the
central obstruction: for continuous variables BP messages are functions,
so the number of message updates is linear in $K$ but the cost of an
update is unbounded unless the messages close in a finite-dimensional
family. This bottleneck fixed the final protocol: prove closure in the
Gaussian case, and for Laplace innovations compare a Gaussian-projected
recursion against a validated grid-based reference
(Chapter~\ref{ch:laplace}).

\emph{A discarded conjecture.} An intermediate conjecture, that the
Laplace $K = 2$ posterior Hessian is approximately rank-two, failed
under numerical tests outside a narrow regime and was dropped; it is
recorded as a negative result in Section~\ref{sec:l-open}.
